%% notes/topics/90-crosscutting/discrepancies.tex
%%
%% 用途:
%%   資料間・ファイル間で内容が食い違っている箇所の台帳。
%%
%% 記入の規則（重要）:
%%   - どちらかを正しいものとして採用しない。
%%     双方の記述を並記し、差分を明示するにとどめる。
%%   - 「どうすれば決着するか」を必ず書く。
%%   - 決着したら、決着の根拠と日付を記録したうえで下部へ移す。
%%     このとき、採用しなかった側の記述も削除せず残す。
%%
%% 種別:
%%   (A) 同名ファイルの複数の版で結論が異なるもの
%%   (B) 引き継ぎ資料の記述と、ディスク上のファイルの実態が一致しないもの
%%   (C) 同一の主張が資料ごとに異なる地位を持つもの
%%   (D) 素朴に統合すると矛盾に見えるが、対象が異なるもの

\section{横断: 未決着の不一致}

\subsection{(A) 原稿の版の食い違い}

\subsection{(B) 改訂履歴の記述とファイルの実態の食い違い}

\subsection{(C) 地位の食い違い}
%% 矛盾ではなく進捗の差。昇格の可否は照合が済むまで判断しない。

\subsection{(D) 見かけの緊張}
%% 矛盾ではないが、統合時に混同しやすいもの。

\subsection{欠落している資料}

\subsection{決着済み}
%% 決着の根拠と日付。採用しなかった側の記述も残す。
